% The back matter: Declarations (with the journal's headings, before the reference list) and the Open Practices
% Statement (immediately prior to the references), as outline.md and the journal's instructions require. Values that
% are known only after the deposit or that the authors give are marked with \tbd.

\backmatter

\section*{Declarations}

\bmhead{Funding}
This research received no specific grant from any funding agency in the public, commercial or not-for-profit sectors.

\bmhead{Conflicts of interest/Competing interests}
The author has no competing interests to declare that are relevant to the content of this article.

\bmhead{Ethics approval}
Not applicable: no participants were recruited. The key presses of every run of a plan were generated by a program;
a person's pointer click is reported in one run (Section~\ref{sec:results-attempts}), and the observation log of the
idle condition of question~5, which holds the key presses that started its sampling, is not deposited. The logs were
checked for key presses that the program did not send, by the stopping rule right after each run and again at deposit;
the blind spots of these checks are stated in Section~\ref{sec:limitations}.

\bmhead{Consent to participate}
Not applicable.

\bmhead{Consent for publication}
Not applicable.

\bmhead{Availability of data and materials}
The dataset described in Section~\ref{sec:dataset}, with its datasheet and the specification of its formats, is
deposited in Zenodo under DOI 10.5281/zenodo.23183438 \citep{dataset}. The two registrations are
on OSF Registries \citep{prereg-macos,prereg-windows}. The plans, the procedures, the phrase corpus, the settings of the
generator and the seed are part of the dataset and of the code repository. Not deposited are the observation log of the
idle condition of question~5, the Windows attempt stopped inside a trial, one aborted attempt of the Windows pilot, and
the Windows pilot runs and an aborted attempt with the second reference field, for the reasons given in
Section~\ref{sec:dataset-excluded}.

\bmhead{Code availability}
The logger, the kit that generates the plans, runs the cells and checks the logs, the analysis code, and the script that
generates the tables, the values quoted in the text and Figures 2 to 5 of this article from the result files are
released in a public repository, \url{https://github.com/mikotomiura/ime-conversion-logging}, with an archived release
under DOI 10.5281/zenodo.23183441 \citep{release}. The analyses were run from exports of the freezing tags of the
registration (Table~\ref{tab:freezes}), with the files that deviations D8 and D10 replace; the supplement lists the
code of every result file. The exploratory analysis of deviation D13 was run with a script that release 1.0.2 adds,
under DOI 10.5281/zenodo.23273986 \citep{release102}; the script imports \texttt{kit/e2/closed\_set.py} of release
1.0.0 unchanged, and the values and the supplementary table of D13 are generated by the version of
\texttt{manuscript/tools/results\_tables.py} in release 1.0.2.

\bmhead{AI declaration}
The author used two tools based on large language models. Claude (Anthropic; the Opus 4.6 to 5.5, Sonnet 4.6 to 5.5
and Haiku 4.5 models, used through the Claude Code command-line tool between February and October 2026) was used to
write all the code of the logger and of the kit (among them the plan generator, the runner, the reference input fields,
the log formats and their validator, the corpus tools and the analysis), the script that generates the tables, the
values quoted in the text and Figures 2 to 5 from the result files, and the code of Figure 1; to draft the two
registrations, the record of deviations, the procedures, the outline of this article and its sections, and the
datasheet of the dataset;
to search for and summarise the documentation and the published work examined; to simulate peer review of the
registrations, of the outline and of this article (the simulated review of this article prompted deviations D11 and
D12, and the assistant, reading the code of the closed-set attack for another study, found the difference between its
description and the code that deviation D13 corrects); to start and monitor the measurement runs on the measuring
machines, and to run the checks of the runs, the packing of the deposits and the analyses, under the author's
instructions (in the first Windows attempt the assistant ended the shell that ran the runner when the machine was short
of memory; Section~\ref{sec:results-attempts}); and to review the code and the text. Codex (OpenAI, used through its
command-line tool) was used to review the code and the text and to give a second opinion on the plan of the
discussion. The research questions are the author's. Parts of the design, of the analysis settings and of the text
were proposed by Claude; the author chose among the options presented, checked and approved them, and froze the
registrations. The text of the results, the discussion, the recommendations and the conclusions was drafted by Claude
from the registered analyses and from the direction and the interpretation that the author gave; the author checked
the drafts, including the changes made after the reviews, and approved them. No language model produced or altered
any value of the data: the plans were generated by the plan generator from the seed, the logs were written by the
programs during the runs, the result values come from the logs through the analysis code (the registered analysis,
the deviations listed in Table~\ref{tab:deviations} and the exploratory analysis) and are copied into the text, or
computed from the result files, by the script, and the figures of the records, such as
the checks of the runs, are copied from the run records. The code was checked by its tests and by the reviews named
above. The interpretation of the results is the author's, and the author is accountable for the content of this
article.

\bmhead{Authors' contributions}
Mikoto Miura is the sole author: conceptualization, methodology, software, validation, investigation, formal analysis,
resources, data curation, visualization, project administration, and writing (original draft, review and editing).

\section*{Open Practices Statement}

The study was pre-registered on OSF Registries in two parts, for the macOS cells (\url{https://osf.io/9nh7r},
DOI 10.17605/OSF.IO/9NH7R) and, in an addendum frozen before the first Windows measurement, for the Windows cell
(\url{https://osf.io/4jwzn}, DOI 10.17605/OSF.IO/4JWZN). The code, the materials and the analysis scripts are available
in a public repository with an archived release (DOI 10.5281/zenodo.23183441), and the data are available under their
own persistent identifier (DOI 10.5281/zenodo.23183438). The analyses were run from exports of the freezing tags, with
the files that deviations D8 and D10 replace, and the exploratory analysis of deviation D13 with a script that
release 1.0.2 adds (DOI 10.5281/zenodo.23273986) and that imports \texttt{kit/e2/closed\_set.py} unchanged; the
deviations from
the registrations are listed in
Table~\ref{tab:deviations}. Not registered are: the analysis of observability (Section~\ref{sec:observability}); the
descriptive statements about what can be computed from the logs (Table~\ref{tab:measures}) and about the uses that the
records support (Section~\ref{sec:dataset-uses}); the recommendations of Section~\ref{sec:recommendations} as
recommendations; the rule by which the abstract and the conclusions select their estimates beyond section~6.1 of the
registration; the exploratory analysis of the distribution of the duration error (Section~\ref{sec:method-analysis});
the exploratory closed-set analysis of deviation D13 (Section~\ref{sec:recoverability-method}); and the comparisons
between the cells or between speeds that the discussion reads. Added after the results had been
seen are the statements of the scenarios that set the ends of the ranges in the abstract and the conclusions, the
pooled values in the conclusions, and the accuracies of question~4 at $K = 10$ and $K = 100$ in recommendation (c) and
the conclusions. Deviations D11 and D12 were made after
the results had been seen, prompted by a simulated review of this article before its submission (AI declaration).
Deviation D13 was also made after the results had been seen, when the code of the closed-set attack was read for
another study; its two feature sets, its rule and its place in this article were fixed before its analysis was run.
The outline that fixed the rules for reporting before the results were analysed is in the public repository as commit
\texttt{cb29a1a9086e77c89cb66c05e39ba4fc8091e074}.
